\documentclass[11pt]{article}
\usepackage[margin=2.5cm]{geometry}
\usepackage{graphicx}
\usepackage{booktabs}
\usepackage{longtable}
\usepackage{array}
\usepackage{amsmath}
\usepackage{amssymb}
\usepackage{siunitx}
\usepackage[colorlinks=true,linkcolor=blue,citecolor=blue,urlcolor=blue]{hyperref}
\usepackage{xcolor}
\usepackage[spanish]{babel}
\decimalpoint % mantener 0.4 en vez de 0,4 (default de babel-spanish)
\usepackage[utf8]{inputenc}
\usepackage{enumitem}
\usepackage{fancyvrb}

\newcolumntype{L}[1]{>{\raggedright\arraybackslash}p{#1}}

\title{Instalación desde cero de la cadena completa\\
MUSIC $\to$ IP-Glasma/3dMCGlauber $\to$ iS3D $\to$ UrQMD $\to$ CRAB3/\texttt{be}\\
\large Checklist con bugs corregidos, para un checkout nuevo en \texttt{xook}}
\author{Bi+Bi correlations campaign}
\date{2026-09-30}

\begin{document}
\maketitle

\section{Alcance}

Esta guía consolida en una sola lista, en orden de ejecución, todos los
bugs de instalación/ejecución ya encontrados y corregidos a lo largo de las
campañas Bi+Bi 5.8/7.7/9~GeV y AuAu 130/200~GeV (fuente:
\texttt{music-install/README.md} \S8, y los README de
\texttt{music/BiBi\_cascade\_*GeV/}), pensada para levantar la cadena
completa en una máquina nueva (\texttt{xook}) sin redescubrir cada uno por
separado. Cierra con un chequeo de sanidad end-to-end: reproducir el
\emph{yield} de referencia Bi+Bi $\sqrt{s_{NN}}=5.8$~GeV medido con una
cascada UrQMD standalone, que es el número contra el que se calibra toda la
cadena hidro en esa energía.

No repite el contenido físico/metodológico de los status reports ya
existentes (\texttt{status\_report\_bibi58\_cascade.tex},
\texttt{status\_urqmd\_crab3.tex}) — esta es una guía operativa de
instalación.

\section{Prerrequisitos del sistema}

\begin{longtable}{L{4cm}L{5cm}L{6cm}}
\toprule
Paquete & Para qué & Instalación (Ubuntu/Debian) \\
\midrule
\endhead
\texttt{g++} / \texttt{gfortran} & compilar MUSIC (C++) y UrQMD (Fortran legacy) & incluidos en \texttt{build-essential} + \texttt{gfortran} \\
\texttt{cmake} & build de MUSIC & \texttt{sudo apt install cmake} \\
\texttt{git} & clonar los repos & \texttt{sudo apt install git} \\
\texttt{libgsl-dev} & interpolación (MUSIC) & \texttt{sudo apt install libgsl-dev} \\
Python 3 + \texttt{numpy}/\texttt{numba}/\texttt{matplotlib} & yields, \texttt{be/} (C2/C3), plots & \texttt{sudo apt install python3 python3-numpy python3-matplotlib}, \texttt{pip install numba} \\
\bottomrule
\end{longtable}

\textbf{Ojo con \texttt{libgsl-dev} específicamente}: si solo está la
librería en tiempo de ejecución (\texttt{libgsl28}) pero no los headers,
\texttt{cmake} compila MUSIC igual pero \emph{sin GSL y sin avisar con un
error fuerte} — el binario queda incompleto y falla de formas confusas más
adelante. Verificar esto explícitamente antes de compilar (o dejar que
\texttt{install\_music.sh} lo haga).

\section{MUSIC (core)}

\begin{enumerate}[leftmargin=*]
  \item Clonar e instalar:
\begin{Verbatim}[fontsize=\small]
git clone https://github.com/isadoji/music-install.git
cd music-install
./install_music.sh /software/mexnica/MUSIC
export MUSIC_DIR=/software/mexnica/MUSIC # agregar a .bashrc
\end{Verbatim}
  \item \textbf{Aplicar parches conocidos antes de dar por buena la
  instalación} 
  \begin{Verbatim}[fontsize=\small]
  ./apply_music_fixes.sh /software/mexnica/MUSIC
\end{Verbatim}
  (\texttt{patches/},este repo):
  \begin{itemize}
    \item \texttt{freeze\_pseudo.patch} --- evita que \texttt{free(particleList)}
      cuelgue en corridas grandes de \texttt{freeze\_pseudo.cpp} (bug de
      limpieza de memoria de MUSIC upstream al final del freeze-out).
    \item \texttt{read\_in\_parameters.patch} --- necesario si se va a usar
      \texttt{EOS\_to\_use 20} (2DTExS): \\el \texttt{read\_in\_parameters.cpp}
      de MUSIC stock solo acepta EOS $\le 14$ de fábrica y rechaza el input
      sin este parche.
  \end{itemize}
  \item \textbf{2DTExS (\texttt{EOS\_to\_use 20})} --- EOS custom, ya
  vendorizada en este repo (\texttt{patches/2dtexs/}): archivos nuevos
  \texttt{eos\_2dtexs.h}/\texttt{eos\_2dtexs.cpp}, los reemplazos completos
  \texttt{eos.cpp}/\texttt{CMakeLists.txt} (los stock de MUSIC no registran
  la EOS 20) y la tabla \texttt{EoS2DTExS.dat} (41~MB, checksum verificado
  contra el original de \texttt{epos/src/MSt/}). Ya no depende de un
  checkout aparte de \texttt{epos/} al lado. Instalar con:
\begin{Verbatim}[fontsize=\small]
./apply_2dtexs.sh /software/mexnica/MUSIC
\end{Verbatim}
  Este script hace todo en un solo paso: copia los archivos nuevos, reemplaza
  \texttt{eos.cpp}/\texttt{CMakeLists.txt}, aplica
  \texttt{freeze\_pseudo.patch} y \texttt{read\_in\_parameters.patch} (si ya
  se corrió \texttt{apply\_music\_fixes.sh} antes, los detecta ya aplicados
  y no falla), copia la tabla a \texttt{\$MUSIC\_DIR/EOS/2DTExS/} y
  recompila. Si solo se necesita la física estándar (\texttt{EOS\_to\_use 9},
  hotQCD) no hace falta correr este paso --- basta con
  \texttt{apply\_music\_fixes.sh} del punto anterior.
  \item Gotchas menores ya manejados automáticamente por los scripts de
  \texttt{music-install} (no requieren acción, se listan para referencia):
  directorio de \texttt{--output} de SLURM debe existir antes de
  \texttt{sbatch} (\texttt{mkdir -p runs/<nombre>}); \texttt{min\_pt 0.01}
  en vez de \texttt{0} exacto (evita el crash de GSL
  \texttt{x values must be strictly increasing}); \texttt{./outputs} se
  crea antes de modo~3; symlinks de \texttt{EOS/}/\texttt{tables/} en el
  directorio de la corrida.
\end{enumerate}

\section{Condiciones iniciales --- IP-Glasma / 3dMCGlauber}

Para Bi+Bi a energías bajas (NICA/BES: 5.8, 7.7, 9~GeV) usar
\textbf{3dMCGlauber} (\texttt{Initial\_profile 13}), no IP-Glasma (que no
tiene descripción nativa de bariones netos y no aplica a estas energías).

\begin{enumerate}[leftmargin=*]
  \item Instalar y parchar el núcleo de Bi-209 (no viene de fábrica):
\begin{Verbatim}[fontsize=\small]
./install_3dmcglauber.sh /software/mexnica/MUSIC/3dMCGlauber
                            # aplica patches/3dmcglauber_add_bi_nucleus.patch
export GLAUBER_DIR=/software/mexnica/MUSIC/3dMCGlauber
\end{Verbatim}
  El parche agrega la parametrización de Bi-209 ($R=6.96$~fm, $a=0.537$~fm,
  convención arXiv:2401.00619).
  \item \textbf{Recortar la salida a 25 columnas} antes de pasarla a MUSIC:
  la versión instalada de 3dMCGlauber escribe 30 columnas
  (\texttt{strings\_event\_N.dat}, incluye campos de transporte de carga
  eléctrica que MUSIC \texttt{public\_stable} todavía no lee). Sin el
  recorte, MUSIC aborta con \texttt{the format of file...is wrong}. Usar
  \texttt{generate\_bibi\_ic.sh} (ya hace el recorte), no la salida cruda de
  \texttt{3dMCGlb.e} directamente.
  \item \texttt{Include\_Rhob\_Yes\_1\_No\_0 1} es \textbf{obligatorio} a
  estas energías (frenado bariónico grande, $\mu_B$ finito), junto con una
  EOS que dependa de $\mu_B$ (\texttt{EOS\_to\_use 14}, \texttt{neos\_bqs},
  o \texttt{EOS\_to\_use 20} si ya se aplicó el parche de 2DTExS de la
  sección anterior) --- nunca combinar el flag con una EOS a $\mu_B=0$
  (\texttt{EOS\_to\_use 2}/hotQCD), sería inconsistente.
  \item Si se usa IP-Glasma en algún punto de comparación (p.\ ej.\ el
  benchmark a 200~GeV): confirmar que existe
  \texttt{nucleusConfigurations/<núcleo>.bin.in} \emph{antes} de correr ---
  su ausencia no da un error claro, produce un \textbf{loop infinito} que
  parece un OOM (consume 24--64--195~GB sin terminar). Con el archivo
  presente el uso real es $\sim$1.2~GB, $\sim$8~min. Además fijar
  \texttt{maxtime 0.4} explícitamente (\texttt{0.6} genera
  \texttt{epsilon-u-Hydro-t0.6-0.dat}, no el \texttt{t0.4} que esperan los
  scripts downstream); builds recientes además renombraron la salida a
  \texttt{epsilon-u-Hydro-TauHydro-0.dat}.
\end{enumerate}

\section{iS3D --- muestreo de partículas sobre la superficie de freeze-out}

\begin{enumerate}[leftmargin=*]
  \item Clonar, parchar y compilar con
\begin{Verbatim}[fontsize=\small]
./install_is3d.sh /software/mexnica/MUSIC/iS3D
\end{Verbatim}
  (este repo) --- ya aplica el parche crítico de abajo automáticamente.
  \item \textbf{Aplicar el parche crítico de \texttt{deltafReader.cpp} antes
  de usarlo para física con bariones.} Bug: \texttt{Deltaf\_Data::calculate\_
  bilinear} indexa la tabla $\delta f$ como \texttt{f\_data[iT][imuB]}, pero
  el loader la guarda como \texttt{f\_data[iB][iT]} (ejes transpuestos). Con
  \texttt{include\_baryon 1}, \emph{todo} coeficiente $\delta f$ (F, G,
  $\beta_{\rm bulk}$, $\beta_V$, $\beta_\pi$) se lee de la celda equivocada
  ---puede además leer fuera de rango si \texttt{iT>80}. Síntoma:
  $\bar p/p$, $\bar n/n$, $\bar\Lambda/\Lambda$ muy asimétricos incluso a
  $\mu_B\approx0$ ($\bar p/p\sim2.4$ en vez de $\sim1$), K/$\pi$ inflado,
  bariones sistemáticamente bajos. El camino \texttt{include\_baryon 0}
  (spline cúbico) \textbf{no} está afectado.
  \begin{itemize}
    \item \textbf{Fix}: parchar \texttt{deltafReader.cpp} para indexar
    \texttt{[iB][iT]} correctamente y recompilar ---
    \texttt{patches/iS3D\_deltafReader\_bilinear\_fix.patch} en este repo,
    aplicado automáticamente por \texttt{install\_is3d.sh}. \textbf{Verificar
    explícitamente que el checkout nuevo en \texttt{xook} lo traiga}, no
    asumirlo por venir del mismo repo.
    \item Workaround si no se puede recompilar de inmediato:
    \texttt{include\_baryon 0} en \texttt{iS3D\_parameters.dat} (fuerza el
    camino spline) --- solo válido si la superficie tiene $\mu_B\approx0$.
  \end{itemize}
  \item Si la superficie viene de una corrida \textbf{boost-invariant}
  (2D), fijar \texttt{mode 8} explícitamente en la config de iS3D --- el
  reader por defecto da resultados raros con superficies 2D. Confirmar con
  $\langle T\rangle$/$\langle\mu_B\rangle$ físicamente razonables tras
  correr.
\end{enumerate}

\section{UrQMD --- afterburner / cascada}

\begin{enumerate}[leftmargin=*]
  \item \textbf{Parche de compilación obligatorio con gfortran $\ge$10}:
  UrQMD 3.4 es Fortran legacy; gfortran $\ge$10 endureció el chequeo de
  tipos por defecto y falla con \texttt{Type mismatch}/\texttt{Rank
  mismatch}.
\begin{Verbatim}[fontsize=\small]
./install_urqmd.sh /software/mexnica/MUSIC/urqmd-3.4
\end{Verbatim}
  (este repo) descarga el tar oficial de \texttt{urqmd.org} y aplica
  \texttt{patches/urqmd\_Linux.mk.patch} automáticamente --- agrega
  \texttt{-std=legacy -fallow-argument-mismatch -ffixed-line-length-none}
  antes de compilar.
  \item \textbf{Semántica invertida de las banderas de salida}: listar una
  unidad de salida (\texttt{fXX}) en la tarjeta de entrada de UrQMD
  \emph{suprime} esa salida (\texttt{bfXX=.true.} $\to$ no escribe), no la
  habilita. Para quedarse solo con \texttt{f19} (OSCAR1997A nativo, el
  formato correcto --- \textbf{no} \texttt{f14}, que es el formato extendido
  propio de UrQMD y \emph{no} es OSCAR), listar las que \emph{no} se
  quieren: \texttt{f13 f14 f15 f16 f18 f20}, dejando \texttt{f19} sin listar.
  \item Si se usa el wrapper del afterburner (no la cascada standalone):
  confirmar que siembra \texttt{frrx/frry/frrz/frr0} (posiciones reales de
  freeze-out) en el loop de \texttt{getoldevent} (\texttt{output.f}) ---
  un wrapper viejo sin este parche deja esas posiciones en cero/no
  inicializadas, invalidando cualquier análisis de $c\tau$/core-halo. El
  parche está en
  \texttt{AuAu\_200GeV/crab\_analysis/patches/urqmd\_afterburner\_frzout\_
  creation\_point.patch}; \textbf{no} se retro-aplica a output ya generado
  antes del fix.
  \item Usar siempre \texttt{f19}/\texttt{f20} nativos para cualquier
  análisis que dependa de la posición real de freeze-out --- \texttt{f20}
  (\texttt{osc99\_event}) escribe \texttt{frrx/frry/frrz/frr0} (coordenadas
  reales); \texttt{f14} (\texttt{file14out}) escribe \texttt{r0/rx/ry/rz}
  (posiciones \emph{free-streamed}, distinta cantidad física) --- confirmado
  leyendo \texttt{output.f} directamente.
  \item \textbf{\texttt{stop 137} no es OOM}: UrQMD 3.4 puede abortar con
  \texttt{stop 137} desde \texttt{anndec.f:285} (canal de decaimiento sin
  probabilidades de rama, p.\ ej.\ \texttt{ityp 139}, $m=2.112$~GeV) --- es
  un exit code propio de UrQMD, no memoria. Si un job array pierde
  3--43\% de tareas según config, revisar primero este crash antes de
  asumir OOM. Mitigación actual: correr en chunks pequeños (no hay
  reintento/skip automático implementado todavía).
\end{enumerate}

\section{CRAB3}

\begin{enumerate}[leftmargin=*]
  \item Clonar, parchar y compilar con
\begin{Verbatim}[fontsize=\small]
./install_crab3.sh /software/mexnica/MUSIC/CRAB3
\end{Verbatim}
  (este repo; CRAB3 no usa \texttt{cmake}, se compila directo con
  \texttt{g++ crab.cpp -o crab.e}). Aplica dos parches:
  \begin{itemize}
    \item \textbf{\texttt{Increase NPHASEMAX in crab.cpp}}: el
    \texttt{\#define NPHASEMAX 20000000} de \texttt{crab.cpp} es un arreglo
    \emph{estático} de punteros (\texttt{NBMAX=1}) que guarda todos los
    $\pi^+$ de todos los eventos pooled del run; campañas grandes
    ($\sim$1M eventos) lo superan y \texttt{crab\_prinput.cpp} aborta. Fix
    vendorizado aquí: \texttt{patches/crab3\_increase\_nphasemax.patch}
    simplemente sube la constante a \texttt{200000000} --- al ser un
    arreglo de punteros (no de objetos), el costo real de memoria (BSS)
    crece solo con las entradas realmente usadas, no con el tamaño
    declarado.
    \item \texttt{patches/crab3\_nmax\_mixing\_heap\_fix.patch}: mueve los
    arreglos de \texttt{MIXED\_PAIRS\_FOR\_DENOM}
    (\texttt{NMAX\_FOR\_MIXING=200000}) de la pila al heap en
    \texttt{crab\_main.cpp} --- evita overflow de stack.
  \end{itemize}
  \item Crear \texttt{results/} antes de correr \texttt{crab.e} ---
  \texttt{fopen()} sin ese directorio segfaultea al arrancar.
  \item \textbf{\texttt{--mem} de SLURM no se puede fijar estático}: la
  multiplicidad real ($\pi^+$/evento) solo se conoce corriendo el
  afterburner primero. Insertar un job intermedio ``autosize'' (muestrea un
  shard real ya escrito, proyecta el total, aplica margen $\sim$30\%) entre
  el afterburner y el job real de CRAB3 --- patrón ya implementado en
  \texttt{music/BiBi\_cascade\_*GeV/figA3/slurm\_*\_crab3\_autosize.sh}.
  \item Si se actualiza memoria de un job ya enviado con
  \texttt{scontrol update JobId=<id> MinMemoryNode=<valor>}: el campo
  espera un entero plano en MB, \emph{no} un sufijo \texttt{G}
  (\texttt{MinMemoryNode=32768}, no \texttt{32G}).
  \item Antes de comparar contra un \texttt{crab3\_qinv*.dat} de
  referencia: leer siempre el header --- algunos usan
  \texttt{k=(p2-p1)/2:} (momento reducido) y otros \texttt{k=(p2-p1):}
  (Qinv completo); \textbf{no} es una convención fija por proyecto, cambia
  por build/config. Un factor $\sim$2 de diferencia en $R$ ajustado suele
  venir de esto.
\end{enumerate}

\section{\texttt{be/} --- C2/C3 en Python}

\begin{enumerate}[leftmargin=*]
  \item \textbf{Unidades}: \texttt{be/kernels.py} trabaja en GeV/fm con
  momento reducido $k=Q_{\rm inv}/2$. Los wrappers de alto nivel
  (\texttt{compute\_c2\_crab}, \texttt{compute\_c2\_batched}) convierten
  \texttt{maxmom\_mev/1000} internamente --- un driver que llama al kernel
  \texttt{numba} directo \textbf{no} hace esa conversión, y el fallo es
  silencioso (exit 0, arrays con la forma correcta, contenido mal: eje
  0--50~GeV en vez de 0--50~MeV, $C_2\approx1.000$ plano). Convertir
  siempre \texttt{maxmom\_gev = maxmom\_mev / 1000.0} antes de binear si se
  llama \texttt{kernels.py} directo; validar contra el estimador
  pool-random como cross-check.
  \item \textbf{Simetrización de Bose--Einstein en C3}: \texttt{weight =
  w12*w13*w23} (producto de pesos pairwise) \textbf{no} es la
  simetrización genuina para $N=3$ bosones idénticos --- hay que sumar
  sobre las $3!=6$ permutaciones de $S_3$. Usar
  \texttt{kernels.triplet\_be\_weight()} (ya suma las 6 fases de
  permutación directo de los cuadrivectores).
  \item \textbf{Memoria}: la implementación original de \texttt{c2.py}/
  \texttt{c3.py} retenía cada par/triplet muestreado en memoria (RSS
  escala linealmente: 100M triplets $\to$ 24~GB, 400M $\to$ 88~GB;
  \texttt{n\_pairs=2e8} $\to$ $\sim$30~GB). Usar el refactor de reservoir
  acotado (\texttt{reservoir\_size}, ya presente) --- RSS acotado
  independiente de \texttt{n\_pairs}/\texttt{n\_triplets}; correr además
  bajo \texttt{ulimit -v} como cinturón de seguridad en jobs grandes.
  \item El reservoir del denominador mixto en
  \texttt{compute\_c3\_diagonal\_crab} debe llenarse \textbf{sin}
  condición desde todo intento crudo, no solo con triplets ya aceptados
  por el corte diagonal --- llenarlo condicionado sesga la referencia
  mixta hacia abajo (meseta sospechosa $\sim$0.5--0.6 a $k$ grande).
  \item \texttt{plotting.py}'s \texttt{plot\_c2()}: enmascarar
  \texttt{ref\_err >= 0} antes de graficar (algunos \texttt{.dat} usan el
  sentinel $-1$ para ``sin datos'' en un bin; sin la máscara,
  \texttt{yerr} negativo crashea el errorbar).
\end{enumerate}

\section{Gotchas generales de SLURM}

\begin{itemize}[leftmargin=*]
  \item \textbf{Nunca fijar \texttt{--time}} en estos jobs --- dejar el
  default del cluster (sin límite); un \texttt{--time} viejo ya causó
  pérdida de 20/60 tareas por \texttt{TIMEOUT} en la campaña de cascada
  UrQMD standalone (\S\ref{sec:validacion}).
  \item \textbf{\texttt{\$0}/\texttt{dirname "\$0"} no localiza el propio
  script bajo un job array}: SLURM copia el script a un spool dir
  (\texttt{/var/lib/slurm/slurmd/jobNNN/}) y ejecuta la copia --- resolver
  rutas relativas a scripts ``vecinos'' desde \texttt{\$SLURM\_SUBMIT\_DIR}
  o hardcodear la ruta absoluta, no derivarla en runtime.
  \item Un \texttt{glob} de muchos archivos de shards falla con
  \texttt{Argument list too long} (excede \texttt{ARG\_MAX}) --- usar
  \texttt{find ... -exec} o iterar desde una lista en archivo.
  \item \texttt{acaxees02} no monta \texttt{/storage} de forma consistente
  --- confirmar el mount antes de lanzar un array que dependa de él.
  \item Los nodos del cluster pueden tener velocidad desigual --- no asumir
  que todas las tareas de un array terminan en tiempo similar.
\end{itemize}

\section{Configurar acceso sin password (para lanzar/monitorear sin dejar
la terminal abierta)}

Si \texttt{xook} necesita transferir datos hacia/desde otro nodo del
cluster por \texttt{rsync}/\texttt{ssh} de forma no interactiva (jobs largos
en \texttt{screen}/\texttt{nohup}), generar una llave dedicada en vez de
depender de password:

\begin{Verbatim}[fontsize=\small]
ssh-keygen -t ed25519 -f ~/.ssh/xook_<destino> -N ""
ssh-copy-id -i ~/.ssh/xook_<destino>.pub usuario@destino
\end{Verbatim}

y usarla explícitamente en el \texttt{rsync}: \texttt{-e "ssh -i
~/.ssh/xook\_<destino>"}. Lanzar transferencias largas con
\texttt{screen -dmS <nombre> bash script.sh} para poder desengancharse
(\texttt{Ctrl+A} luego \texttt{D}, \textbf{no} \texttt{Ctrl+D}) sin matar el
proceso.

\section{Validación: yield Bi+Bi $\sqrt{s_{NN}}=5.8$~GeV vs.\ cascada UrQMD standalone}
\label{sec:validacion}

Antes de confiar en la cadena hidro completa (MUSIC~$\to$~iS3D~$\to$~UrQMD~
$\to$~CRAB3) recién instalada en \texttt{xook}, se recomienda reproducir el
número de referencia contra el que se calibra \texttt{s\_factor}/los
parámetros de \texttt{3dMCGlauber} a esta energía: el yield final de
$\pi^+$ de una cascada UrQMD \textbf{standalone} (sin MUSIC, sin iS3D ---
solo UrQMD puro), que sirve de ``blanco'' independiente de cualquier
suposición hidrodinámica.

\subsection{Instalación de la cascada standalone}

Binario distinto al afterburner del pipeline hidro (\texttt{urqmd.x86\_64}
vs.\ \texttt{urqmd\_afterburner}) --- mismo repo \texttt{urqmd-3.4}, mismo
parche de compilación de la \S6 (\texttt{Linux.mk.patch}, obligatorio con
gfortran $\ge$10), pero compilado en modo standalone (genera su propia
colisión núcleo-núcleo en vez de leer una superficie de freeze-out de
iS3D). Instalar con \texttt{./install\_urqmd.sh} (este repo).

\subsection{Tarjeta de entrada}

Reconstruida leyendo \texttt{input.f} (sin documentación externa
disponible):

\begin{Verbatim}[fontsize=\small]
pro 209 83        # proyectil: Bi (A=209, Z=83)
tar 209 83        # blanco:    Bi (A=209, Z=83)

nev <N>           # numero de eventos
IMP 0.0 1.0       # parametro de impacto explicito (fm), central

ecm 5.8           # sqrt(s_NN) por par NN (srtflag=1)
eos 0             # 0 = cascada pura, sin potenciales
rsd <seed>
tim 200 200       # tiempo total de propagacion / intervalo de salida (fm/c)

f13
f14
#f19              # NO listar: es la salida OSCAR1997A real que se necesita
f15
f16
f20               # listar = suprimir esa unidad (ver bug de semantica, S5)

xxx               # fin de tarjeta
\end{Verbatim}

Ejecución: variables de entorno \texttt{ftn09}=archivo de entrada,
\texttt{ftn13..ftn20}=archivos de salida por unidad Fortran, luego
\texttt{./urqmd.x86\_64}.

\textbf{Gotcha ya encontrado en esta misma validación}: el piloto original
se corrió creyendo que \texttt{f14} era el formato OSCAR1997A --- es
incorrecto (\texttt{f14} es el formato extendido nativo de UrQMD, columnas
\texttt{t X Y Z E Px Py Pz m ityp iso chg lcl\# ncl or}). \texttt{f19} es
el OSCAR1997A verdadero (con PDG ids vía \texttt{pdgid()}). Confirmar que
la tarjeta de entrada deja \texttt{f19} habilitado (sin listar) antes de
lanzar la campaña completa, no solo un evento de prueba.

\subsection{Paralelizar en SLURM}

Job array, cada tarea en un directorio aislado (\texttt{mktemp -d}) para
evitar condiciones de carrera al escribir archivos de nombre fijo:

\begin{Verbatim}[fontsize=\small]
sbatch --array=0-59 scripts/slurm_urqmd_cascade_array.sh \
    <n_total_events> <outdir> <base_seed>
\end{Verbatim}

Sin \texttt{--time} (\S9) --- un límite de tiempo viejo hizo perder 20/60
tareas por \texttt{TIMEOUT} en la corrida de referencia original.

\subsection{Cálculo del yield y comparación con el blanco}

Sobre el \texttt{f19} combinado (OSCAR1997A, filtrando solo $\pi^+$),
calcular la rapidez real por partícula,
\[
  y = \tfrac{1}{2}\ln\!\left(\frac{E+p_z}{E-p_z}\right),
\]
y contar $\pi^+$ con $|y|<0.5$:
\[
  \frac{dN_{\pi^+}}{dy}\Big|_{|y|<0.5} = \frac{N_{\pi^+}(|y|<0.5)}{N_{\rm eventos}}.
\]

\textbf{Valor de referencia ya medido} (job 12333, 66{,}670 eventos
sobrevivientes de 100{,}000 lanzados --- 20/60 tareas se perdieron por el
\texttt{TIMEOUT} descrito arriba, no relanzadas):
\[
  \frac{dN_{\pi^+}}{dy}\Big|_{|y|<0.5} = \frac{5{,}717{,}085}{66{,}670} = 85.75.
\]

\textbf{Criterio de aceptación para el checkout nuevo en \texttt{xook}}:
correr el mismo número de eventos (o una fracción estadísticamente
suficiente, $\gtrsim$10{,}000) con la misma tarjeta de entrada
(\texttt{ecm 5.8}, \texttt{IMP 0.0 1.0}, \texttt{eos 0}) y confirmar que el
$dN_{\pi^+}/dy(|y|<0.5)$ medido cae dentro del error estadístico de 85.75.
Una discrepancia grande indica que algún parche de la \S6 (semántica de
banderas \texttt{fXX}, o el parche de compilación de gfortran) no se aplicó
correctamente en la instalación nueva, \textbf{antes} de invertir tiempo de
cómputo en la cadena hidro completa.

Este mismo yield (85.75) es además el ``blanco de calibración'' usado
después para ajustar \texttt{3dMCGlauber}/\texttt{s\_factor} en las
corridas hidro completas a esta energía (ver
\texttt{BiBi\_cascade\_5.8GeV/README.md}, secciones de escaneo de
\texttt{s\_factor} y barrido de \texttt{shadowing\_factor}/\texttt{BG}) ---
reproducirlo aquí no es solo un chequeo de instalación, es el primer paso
real de esa calibración.

\section{Resumen: orden de instalación recomendado}

\begin{enumerate}
  \item Prerrequisitos del sistema (\S2).
  \item MUSIC + parches \texttt{freeze\_pseudo}/\texttt{read\_in\_parameters}
  (\S3).
  \item 3dMCGlauber + parche de núcleo Bi-209 (\S4).
  \item UrQMD (standalone \emph{y} afterburner) + parche de compilación
  gfortran (\S6, \S\ref{sec:validacion}).
  \item \textbf{Validación de yield 5.8~GeV vs.\ cascada UrQMD}
  (\S\ref{sec:validacion}) --- antes de seguir.
  \item iS3D + parche crítico \texttt{deltafReader.cpp} (\S5).
  \item CRAB3 (\S7).
  \item \texttt{be/} (\S8), solo si se van a correr correlaciones C2/C3
  además del pipeline de yields/espectros.
\end{enumerate}

\end{document}
